\section{References and Back Matter}

\subsection{Abbreviations}

\vspace{0.30em}

\begin{xltabular}{\textwidth}{L{2.3cm} Y L{2.3cm} Y}
\toprule
\textbf{Term} & \textbf{Meaning} & \textbf{Term} & \textbf{Meaning} \\
\midrule
\endfirsthead
\multicolumn{4}{@{}l}{\footnotesize\itshape Table continued from the previous page.}\\
\toprule
\textbf{Term} & \textbf{Meaning} & \textbf{Term} & \textbf{Meaning} \\
\midrule
\endhead
\midrule
\multicolumn{4}{r@{}}{\footnotesize\itshape Table continued on the next page.}\\
\endfoot
\bottomrule
\endlastfoot
AE & Adverse event & IRB & Institutional Review Board \\
BMI & Body mass index & ISGPS & International Study Group on Pancreatic Surgery \\
CA 19-9 & Carbohydrate antigen 19-9 & ISM & Independent safety monitor \\
CRF & Case report form & LLM & Large language model \\
CRO & Contract research organisation & OS & Overall survival \\
DL1 to DL3 & Dose levels 1 to 3 & PDAC & Pancreatic ductal adenocarcinoma \\
DLT & Dose-limiting toxicity & PFS & Progression-free survival \\
DSMB & Data and Safety Monitoring Board & R0 & Microscopically clear resection margin \\
ECOG & Eastern Cooperative Oncology Group & RECIST & Response Evaluation Criteria in Solid Tumors \\
e-stop & Emergency stop & SAE & Serious adverse event \\
FDA & Food and Drug Administration & SMA & Superior mesenteric artery \\
IDE & Investigational Device Exemption & SMV & Superior mesenteric vein \\
IND & Investigational New Drug & USL & Unified Safety Level \\
\end{xltabular}

\vspace{0.30em}

\subsection{Figure inventory}

Thirty figures, numbered 1 to 30, with no gaps and no duplicates. The counts by
type are Mermaid 9, D2 7, PlantUML 5, Graphviz 5, and Diagrams (Python) 4, which
follows the purpose each type serves rather than an equal quota.

\vspace{0.30em}

\begin{xltabular}{\textwidth}{C{1.0cm} L{3.6cm} C{1.0cm} Y}
\toprule
\textbf{Fig} & \textbf{Type} & \textbf{\S} & \textbf{Subject} \\
\midrule
\endfirsthead
\multicolumn{4}{@{}l}{\footnotesize\itshape Table continued from the previous page.}\\
\toprule
\textbf{Fig} & \textbf{Type} & \textbf{\S} & \textbf{Subject} \\
\midrule
\endhead
\midrule
\multicolumn{4}{r@{}}{\footnotesize\itshape Table continued on the next page.}\\
\endfoot
\bottomrule
\endlastfoot
1 & Mermaid-type & 1 & The seven commitments made to the participant \\
2 & Graphviz-type & 1 & Accountability chain rooted at the participant \\
3 & Mermaid-type & 2 & Journey schema with every decision point marked \\
4 & D2-type & 2 & Schedule of Activities from the participant's chair \\
5 & D2-type & 3 & Twenty-one concerns in six disjoint families \\
6 & Graphviz-type & 3 & Every concern wired to its answering clause \\
7 & Mermaid-type & 3 & Concern prevalence against answer completeness \\
8 & PlantUML-type & 3 & Safeguards as capabilities the participant holds \\
9 & Diagrams (Python) & 3 & Where each concern physically lives \\
10 & Mermaid-type & 4 & Objective, endpoint, and the sentence it authorises \\
11 & D2-type & 4 & Endpoint registry with a plain-meaning field \\
12 & PlantUML-type & 5 & Participant status as a guarded state machine \\
13 & Graphviz-type & 5 & The 3+3 rule as an auditable decision tree \\
14 & Mermaid-type & 6 & Eligibility as a two-way gate \\
15 & PlantUML-type & 6 & Direct booking and sponsor response \\
16 & D2-type & 6 & The five choices the participant keeps \\
17 & Diagrams (Python) & 7 & Operating room and on-premises stack \\
18 & PlantUML-type & 7 & Emergency-stop timing budget \\
19 & Mermaid-type & 7 & Advise, approve, execute, and log \\
20 & D2-type & 7 & Force caps and vascular no-fly envelope \\
21 & Graphviz-type & 7 & Hazard fault tree with named barriers \\
22 & PlantUML-type & 8 & Four withdrawal routes and their data outcomes \\
23 & Mermaid-type & 8 & Consent lifecycle with forced re-consent \\
24 & Mermaid-type & 9 & Participant visit calendar with totals \\
25 & Diagrams (Python) & 9 & Data pipeline with per-store read lists \\
26 & D2-type & 9 & Ten robot instruction cards, PDAC-scoped \\
27 & Mermaid-type & 10 & Five-panel quantitative reassurance dashboard \\
28 & Graphviz-type & 10 & Evidence provenance for every quoted number \\
29 & D2-type & 11 & Responsibility matrix, one accountable per row \\
30 & Diagrams (Python) & 12 & Post-trial continuity and the payer per item \\
\end{xltabular}

\vspace{0.30em}

\subsection{Table inventory}

Forty-three tables. Every one is set to the body text width, every fixed column
is ragged-right, and every table of more than two rows is breakable with a
repeated header, so no table can run off the foot of a page.

\vspace{0.30em}

\begin{xltabular}{\textwidth}{C{1.0cm} L{4.6cm} Y}
\toprule
\textbf{\S} & \textbf{Table} & \textbf{What it carries} \\
\midrule
\endfirsthead
\multicolumn{3}{@{}l}{\footnotesize\itshape Table continued from the previous page.}\\
\toprule
\textbf{\S} & \textbf{Table} & \textbf{What it carries} \\
\midrule
\endhead
\midrule
\multicolumn{3}{r@{}}{\footnotesize\itshape Table continued on the next page.}\\
\endfoot
\bottomrule
\endlastfoot
1 & The seven commitments & Each commitment against the clause that makes it enforceable \\
1 & Diagram types and counts & Which question each of the five types answers \\
1 & How to check a commitment & The artefact behind each commitment, and who holds it \\
2 & Glossary & Twenty terms in the sense this protocol uses them \\
2 & Schedule of Activities & The conventional matrix, retained for a treating oncologist \\
3 & Concern families & Twenty-one concerns grouped and mapped to sections \\
3 & Answered by governance alone & The five, and why no harder answer is available \\
3 & The twenty-one, one row each & Concern, question, answer, and answer class \\
4 & Endpoint table & Every endpoint with its timepoint and plain meaning \\
4 & Endpoint cost & What each endpoint costs the participant to measure \\
4 & Risk classes & Which risks are new to this study and what bounds them \\
5 & Dose levels & What has to be true before each level opens \\
5 & Stopping rules & Seven prespecified triggers, four of them device-specific \\
6 & Two-way gate asymmetry & What differs between the protocol and participant columns \\
6 & Inclusion criteria & Eight criteria, each with the reason it exists \\
6 & Exclusion criteria & Seven criteria, each with the reason it exists \\
6 & Subgroup reporting & What the study can and cannot say at $n \leq 18$ \\
7 & Operative limits & Eight limits with cap, worst observed, and what each protects \\
7 & The day, hour by hour & Each step of the operative day, marked study-specific or not \\
7 & Difference from a conventional robotic Whipple & Ten aspects, four constraints added, nothing changed about the operation \\
8 & Withdrawal routes & Four routes and the data consequence of each \\
8 & Two undrawn transitions & Why re-consent and its refusal matter \\
8 & Study-initiated discontinuation & Six triggers and what follows for the participant \\
9 & Visit burden & Visits, inpatient days, clinic days, and totals \\
9 & Store access and retention & Seven stores, retention in years, and the read list \\
10 & Provenance classes & The four class letters and what to do with each \\
10 & Learning-curve comparator & Fistula and mortality by phase, with the protocol response \\
10 & Safety-limit headroom & Six limits with cap, worst observed, and headroom \\
11 & Responsibility matrix & Nine failure classes, accountable party, who to call \\
11 & Oversight bodies & Independence and the strongest action each may take \\
11 & Breach routes & Three likely breaches and the independent route for each \\
12 & Bill provisions & What H. R. 9510 v5 would establish, against current practice \\
12 & Rights during the study & Eleven rights, marked regulatory minimum or beyond \\
12 & Costs & Eleven items, the payer for each, and the two unsettled \\
13 & Abbreviations, figures, tables, version history & This back matter \\
\end{xltabular}

\vspace{0.30em}

\subsection{Version history}

\vspace{0.30em}

\begin{xltabular}{\textwidth}{L{2.0cm} L{2.6cm} Y}
\toprule
\textbf{Version} & \textbf{Date} & \textbf{Change} \\
\midrule
\endfirsthead
\multicolumn{3}{@{}l}{\footnotesize\itshape Table continued from the previous page.}\\
\toprule
\textbf{Version} & \textbf{Date} & \textbf{Change} \\
\midrule
\endhead
\midrule
\multicolumn{3}{r@{}}{\footnotesize\itshape Table continued on the next page.}\\
\endfoot
\bottomrule
\endlastfoot
Draft 1.0 & July 31, 2026 & First release. Thirteen sections, thirty figures across five diagram types, forty-three tables, twenty-one documented concerns answered clause by clause, and an eight-stage build from machine-readable diagram sources through a bracketed draft and a populated full version to a polished final source set. \\
\end{xltabular}

\vspace{0.30em}

\bmhead{Data and code availability}

Every figure source, every LaTeX source, and every stage narrative is public at
\href{https://github.com/kevinkawchak/robotic-surgeries/tree/main/patient-robot-advocacy}{github.com/kevinkawchak/robotic-surgeries/tree/main/patient-robot-advocacy}.
The five diagram-source directories hold the machine-readable originals in
Mermaid, PlantUML, D2, Python, and DOT; the three paper directories hold the
LaTeX at each stage together with an Overleaf-ready archive. No participant data
exists, because no participant exists: this is an independent advocacy paper
about a protocol, not a report of a conducted study.

\bmhead{Acknowledgments}

The concern evidence in \S 3 rests on two dated research passes of July 28,
2026, one by Gemini 3.1 Pro and one by ChatGPT 5.6 Thinking Extended, filed in
the repository at \texttt{research/}. The parent clinical protocol is the
author's own \cite{onpremwhippl}. The LaTeX build across all eight stages was
adapted using Claude Code Opus 5.

\bmhead{Ethical disclosures}

No human participants, no animals, and no identifiable data. This is an
independent research paper and a practical adoption guide. It is not an active
IND or IDE, not an approved protocol, and not medical or regulatory advice. It
is not endorsed by the FDA, the NIH, HHS, an Institutional Review Board, the
International Council for Harmonisation, or any sponsor, contract research
organisation, site, regulator, or medical society. The author declares no
competing interest other than authorship of the cited works.

\bmhead{Rights and permissions}

Released under CC BY 4.0. Reproduced United States Government regulatory text is
used under 17 U.S.C. \S 105.

\bmhead{Cite this article}

\vspace{0.30em}

\noindent Kawchak K. \textit{Patient Robot Advocacy: A Phase 1, First-in-Human,
PDAC Clinical Trial Protocol of a LLM-Directed Robotic Whipple with Daraxonrasib
(RMC-6236)}. Draft 1.0. Zenodo; 2026.
\href{https://doi.org/10.5281/zenodo.xxxxxxxx}{10.5281/zenodo.xxxxxxxx}.
Repository v1.0.0:
\href{https://github.com/kevinkawchak/robotic-surgeries/tree/main/patient-robot-advocacy}{github.com/kevinkawchak/robotic-surgeries/tree/main/patient-robot-advocacy}.

\bmhead{References}

\bibliographystyle{ieeetr}
\bibliography{references}

\bmhead{How this paper may be reused}

The paper is released under CC BY 4.0, and reproduced United States Government
regulatory text is used under 17 U.S.C. \S 105. Two reuses are anticipated and
are worth naming, because both would be improvements on the paper's current
state.

A site may adapt \S 1, \S 2, and \S 9 into its own participant-facing material,
substituting its own numbers. Nothing in those sections is specific to a single
sponsor, and the artefact list in \S 1.2 and the instruction sheets in \S 9.7
are the parts most directly reusable. The requirement is attribution and the
removal of any number the adapting site cannot itself produce.

A reader may check the paper against its sources. Every figure has a
machine-readable original in one of the five diagram-source directories, every
number in \S 10 carries a provenance letter, and every claim about the parent
protocol resolves to a clause in it. A disagreement found this way is worth
more to the author than agreement, and the escalation route for a factual
correction is the repository's issue tracker rather than the trial's.
